\chapter{Entreprises natives de la crypto}\label{ch:in:crypto-native-firms}

La plus grande plateforme d'échange de crypto-actifs cotée aux États-Unis déclarait 4\,510 salariés à la fin de 2022, après une restructuration menée en juin, qui
avait réduit d'environ 18~\% ses effectifs~; en janvier 2023, elle en a retranché environ 21~\% de plus, et un an plus tard elle comptait 3\,416 salariés. Son chiffre d'affaires avait
chuté de 59~\% en 2022. En 2024, son chiffre d'affaires a plus que doublé, et à la fin de 2025, elle employait 4\,951 personnes. Dans la crypto,
les effectifs suivent le chiffre d'affaires, lui-même tributaire du marché des actifs que le secteur échange, avec un décalage d'environ un an~;
et de nombreux salariés sont aussi payés dans ces mêmes actifs. Ce chapitre décrit les entreprises natives de la crypto en tant qu'employeurs~: les types qui existent, ce qui distingue le travail dans ces entreprises,
comment leurs effectifs évoluent au fil du cycle, ce que vaut une \omterm{def:in:crypto-native-firms:tokens}{rémunération en jetons}, et ce qu'il advient d'un salarié
lorsqu'une telle entreprise fait faillite.

\section{Types d'employeurs}

Les entreprises natives de la crypto sont celles dont l'activité a commencé dans les crypto-actifs plutôt que d'y avoir migré.
\begin{itemize}
\item \textbf{Teneurs de marché et sociétés de trading} cotent sur les plateformes centralisées et négocient on-chain (livre~11,
chapitres~24--25), et font du market making pour des émetteurs de jetons dans le cadre d'accords (livre~3, chapitre~24).
\item \textbf{Fonds} investissent dans des jetons et dans les sociétés et protocoles qui les émettent.
\item \textbf{Bourses et courtiers} exploitent des plateformes centralisées (livre~3, chapitre~15), assurent la conservation et les passerelles entre monnaie et
crypto-actifs.
\item \textbf{Équipes de protocole et fondations} construisent et financent le logiciel d'une blockchain ou d'une application qui s'y exécute.
\end{itemize}

\begin{definition}[Fondation de protocole]\label{def:in:crypto-native-firms:foundation}
Une \emph{fondation de protocole}\index{fondation de protocole} est une organisation à but non lucratif qui finance et coordonne la
développement d'un protocole blockchain à code source ouvert et de son écosystème, généralement à partir d'une trésorerie détenue en partie dans le
jeton propre au protocole, et qui ne possède pas le protocole qu'elle soutient.
\end{definition}

Une fondation se présente comme «~une organisation à but non lucratif qui soutient l'écosystème Ethereum~» et ajoute qu'elle n'est «~pas une
entreprise technologique, ni une association ordinaire à but non lucratif~». Sa politique de trésorerie montre ce que le financement par un jeton signifie pour les personnes qu'elle rémunère~:
«~vendre périodiquement de l'ETH pour garantir des réserves suffisantes pour les années à venir, et augmenter de façon programmée nos réserves en monnaie fiduciaire
pendant les phases haussières pour financer les dépenses pendant les phases baissières~». Le budget d'une fondation, et donc ses recrutements, dépend du
cours des actifs qu'elle détient.

\section{Ce qui distingue~: opérations en continu, conservation, rémunération en jetons, statut réglementaire}

Quatre éléments distinguent les employeurs natifs de la crypto du reste du secteur.
\begin{enumerate}
\item \textbf{Les marchés ne ferment jamais} (livre~3, chapitre~29, appelle cette pratique trading follow-the-sun)~: les astreintes,
les incidents de week-end et les équipes réparties dans le monde sont la norme, et non l'exception (chapitre~26).
\item \textbf{La conservation est un problème propre à l'entreprise.} Les clés, les portefeuilles et la sécurité opérationnelle qui les entourent constituent une
tâche d'ingénierie centrale (livre~3, chapitre~24), et une défaillance est une perte, non une panne.
\item \textbf{La rémunération est souvent en partie versée en jetons}, dont la valeur varie avec le marché et dont le traitement fiscal diffère selon les
pays (section~4).
\item \textbf{Le statut réglementaire reste en cours de définition} dans de nombreuses juridictions (livre~3, chapitre~24, sur le
règlement de l'UE sur les marchés de crypto-actifs (MiCA))~: les entreprises déménagent, obtiennent des agréments et se restructurent à mesure que les règles entrent en vigueur, et les salariés les suivent.
\end{enumerate}

\section{Effectifs au fil du cycle}

\begin{dated}{2025-12}{Une plateforme de crypto-actifs cotée au fil du cycle}\label{dat:in:crypto-native-firms:coinbase}
Coinbase (formulaires 10-K pour 2022 à 2025)~: chiffre d'affaires total de 7\,839~millions de dollars (2021), 3\,194~millions de dollars (2022), 3\,108~millions de dollars
(2023), 6\,564~millions de dollars (2024), 7\,181~millions de dollars (2025)~; effectifs en fin d'année 4\,510 (2022), 3\,416 (2023), 3\,772 (2024),
4\,951 (2025). En juin 2022, une restructuration a touché environ 18~\% des effectifs (environ 1\,100 salariés)~; en janvier 2023, une
deuxième a touché environ 21~\% des effectifs de fin 2022 (environ 950 salariés).
\end{dated}

Le schéma est décalé (la \cref{fig:in:crypto-native-firms:cycle}). Le chiffre d'affaires a chuté de 59~\% en 2022 et de 3~\% en 2023~; les effectifs
ont baissé de 24~\% en 2023, l'année qui suit l'effondrement du chiffre d'affaires. Le chiffre d'affaires a progressé de 111~\% en 2024 et de 9~\% en 2025~; les effectifs ont progressé de 10~\% en
2024 et de 31~\% en 2025. Une entreprise embauche sur la base du chiffre d'affaires de l'année précédente, et réduit aussi ses effectifs sur ce même critère, si bien qu'un salarié qui arrive au sommet d'un cycle
rejoint la cohorte la plus exposée à la prochaine réduction d'effectifs. Cinq années d'une seule entreprise ne font pas une loi, mais c'est le schéma qu'un candidat à un poste dans la crypto
doit anticiper tant que les données ne disent pas le contraire.

\begin{omfigure}
\begin{tikzpicture}
\begin{axis}[width=10.5cm,height=5.4cm,xmin=2020.6,xmax=2025.4,ymin=0,ymax=9,xtick={2021,2022,2023,2024,2025},
  xticklabel style={/pgf/number format/1000 sep={}},ylabel={\small milliards de \$~; salariés (milliers)},
  tick label style={font=\footnotesize},grid=major,grid style={gray!20},table/col sep=comma,unbounded coords=jump,
  legend style={legend cell align=left,font=\scriptsize,at={(0.5,-0.22)},anchor=north,/tikz/every even column/.append style={column sep=8pt}},legend columns=2]
\addplot[omDef,thick,mark=*] table[x=year,y=revenue]{figdata/industry/10-crypto-native-firms/coinbase.csv};
\addplot[omThm,thick,mark=square*] table[x=year,y=employees]{figdata/industry/10-crypto-native-firms/coinbase.csv};
\legend{chiffre d'affaires total (milliards de \$),salariés en fin d'année (milliers)}
\end{axis}
\end{tikzpicture}

\omcaption{Une plateforme de crypto-actifs cotée, 2021--2025~: chiffre d'affaires total et salariés en fin d'année. Les effectifs de 2021 ne figurent pas dans les
documents consultés. Les effectifs suivent le chiffre d'affaires avec un décalage d'environ un an. Données~: Coinbase, formulaires 10-K pour 2022--2025, via
\texttt{in\_crypto.headcount}.}\label{fig:in:crypto-native-firms:cycle}
\end{omfigure}

Le chiffre d'affaires par salarié reproduit le même cycle, vu du côté du salarié (la \cref{fig:in:crypto-native-firms:perhead})~: 0.71
million de dollars par salarié en fin d'année en 2022, 0.91 million de dollars en 2023, 1.74 millions de dollars en 2024 et 1.45 millions de dollars en 2025, un
écart d'un facteur 2.5 en deux ans. Le chapitre~12 a mesuré la même sensibilité d'une entreprise à l'autre~: le revenu par tête dans la crypto
a évolué à contre-courant de la volatilité du marché des actions, et non avec elle. Une enveloppe de bonus (bonus pool) fixée sur le chiffre d'affaires suit cette courbe~; un
effectif fixé d'après le chiffre d'affaires de l'année précédente la suit avec retard, et c'est à ce moment que tombent les réductions d'effectifs.

\begin{omfigure}
\begin{tikzpicture}
\begin{axis}[ybar,width=10.5cm,height=4.8cm,xmin=2021.4,xmax=2025.6,ymin=0,ymax=2,bar width=16pt,
  xtick={2022,2023,2024,2025},xticklabel style={/pgf/number format/1000 sep={}},
  ylabel={\small millions de \$ par salarié},tick label style={font=\footnotesize},grid=major,grid style={gray!20},
  table/col sep=comma,nodes near coords,nodes near coords style={font=\scriptsize},
  nodes near coords={\pgfmathprintnumber[fixed,precision=2]{\pgfplotspointmeta}}]
\addplot[fill=omDef,draw=none] table[x=year,y=perhead]{figdata/industry/10-crypto-native-firms/perhead.csv};
\end{axis}
\end{tikzpicture}

\omcaption{Chiffre d'affaires total par salarié en fin d'année d'une plateforme de crypto-actifs cotée, 2022--2025. Données~: Coinbase, formulaires 10-K,
via \texttt{in\_crypto.per\_head}.}\label{fig:in:crypto-native-firms:perhead}
\end{omfigure}

\section{Rémunération en jetons}

\begin{definition}[Rémunération en jetons, bon de souscription de jetons]\label{def:in:crypto-native-firms:tokens}
La \emph{rémunération en jetons}\index{rémunération en jetons} (token compensation) est une rémunération versée en crypto-actif, généralement selon un calendrier d'acquisition (vesting)
assorti d'une période de carence, et souvent d'une période de blocage après l'acquisition, pendant laquelle les jetons acquis ne peuvent pas être vendus. Un \emph{bon de souscription
de jetons}\index{bon de souscription de jetons} (token warrant) est un droit, accordé avec des titres de capital d'une entreprise, de recevoir des jetons que l'entreprise ou une
société affiliée pourra émettre à l'avenir, en proportion des titres de capital détenus.
\end{definition}

Les autorités fiscales traitent les jetons reçus en contrepartie d'un travail comme une rémunération. Les indications de l'administration fiscale américaine précisent que la monnaie virtuelle versée
par un employeur en rémunération de services constitue un salaire aux fins des cotisations d'emploi~; celles de l'administration britannique indiquent que les crypto-actifs reçus au titre d'un revenu d'emploi
sont considérés comme un avantage en nature et sont soumis à l'impôt sur le revenu et aux cotisations nationales d'assurance sur
leur valeur. La valeur imposée est la valeur à la réception, généralement à l'acquisition. Si le jeton baisse ensuite pendant une période de blocage, le
salarié a payé l'impôt, en numéraire, sur une valeur qui n'existe plus.

\begin{method}[Évaluer une attribution de jetons]\label{met:in:crypto-native-firms:grant}
\begin{enumerate}
\item Établir le calendrier~: la tranche de carence, les tranches mensuelles, le blocage après chaque acquisition.
\item Simuler le prix du jeton dans la devise de l'attribution sur la période d'acquisition et la période de blocage, avec une volatilité et une dérive
énoncées (paramètres, et non prévisions).
\item Pour chaque tranche, calculer l'impôt dû à l'acquisition et le produit à la première date à laquelle elle peut être vendue, net de toute
décote de liquidité.
\item Présenter la distribution des produits nets face à la même valeur d'attribution versée en numéraire, et la fréquence à laquelle l'impôt d'une tranche
dépasse son produit.
\end{enumerate}
\end{method}

\begin{example}[Une attribution sur quatre ans]\label{ex:in:crypto-native-firms:grant}
Une attribution d'une valeur de 400\,000~\$ au cours du jeton à la date d'attribution est acquise à hauteur de 25~\% après une carence d'un an, puis mensuellement sur
trois ans~; chaque tranche est bloquée pendant six mois~; l'impôt représente 40~\% de sa valeur à l'acquisition (illustratif). Avec une
volatilité de 80~\% par an et sans dérive, le résultat net médian sur 20\,000 trajectoires simulées est de 106\,193~\$, contre 240\,000~\$
pour la même attribution versée en numéraire~; les 10\up{e} et 90\up{e}~centiles sont de 17\,467~\$ et 537\,810~\$~; la moyenne, 235\,062~\$, est proche
du montant en numéraire. Dans 76.2~\% des trajectoires, l'impôt d'au moins une tranche dépasse ce qu'elle rapporte ensuite à la vente
(\cref{fig:in:crypto-native-firms:grant}). À une volatilité de 40~\%, la médiane est de 194\,318~\$ et une telle tranche apparaît dans
2.4~\% des trajectoires.
\end{example}

\begin{omfigure}
\begin{tikzpicture}
\begin{axis}[width=10.5cm,height=5.4cm,ybar,bar width=4.5pt,xmin=-100,xmax=1000,ymin=0,
  xlabel={\small résultat net de l'attribution (milliers de \$)},ylabel={\small part des trajectoires (~\%)},
  tick label style={font=\footnotesize},grid=major,grid style={gray!20},table/col sep=comma]
\addplot[fill=omDef!55,draw=omDef] table[x=mid,y=share]{figdata/industry/10-crypto-native-firms/grant_hist.csv};
\draw[omThm,thick,dashed] (axis cs:240,0) -- (axis cs:240,26);
\node[font=\scriptsize,omThm,anchor=south west] at (axis cs:245,20) {équivalent en numéraire 240~k\$};
\end{axis}
\end{tikzpicture}

\omcaption{Le résultat net d'une attribution de jetons de 400\,000~\$ sur quatre ans (acquisition de 25~\% à un an, puis mensuellement,
blocage de six mois, impôt de 40~\% à l'acquisition), 20\,000 trajectoires de prix simulées à une volatilité annuelle de 80~\% et sans dérive~; les résultats au-delà
de 1~million de dollars sont comptés dans la dernière classe. La ligne en pointillés représente la même attribution versée en numéraire. Paramètres illustratifs. Données~:
\texttt{firm.tokencomp.simulate}, via \texttt{in\_crypto.grant\_hist}.}\label{fig:in:crypto-native-firms:grant}
\end{omfigure}

La distribution est asymétrique~: la plupart des trajectoires laissent au salarié moins que ne lui aurait rapporté le numéraire, une minorité beaucoup plus.
La moyenne est proche de la valeur en numéraire parce que le prix n'a pas de dérive, ce qui explique précisément pourquoi la médiane ne l'est pas~: à une volatilité de 80~\%,
la trajectoire typique perd de la valeur. Un candidat qui compare une attribution de jetons à du numéraire compare une loterie
d'espérance à peu près égale à une certitude, avant de compter l'impôt qui porte sur la valeur à l'acquisition.

\section{Quand l'entreprise fait faillite~: le salarié dans le dossier de faillite}

Les entreprises de crypto-actifs peuvent faire faillite rapidement (livre~3, chapitre~15, et livre~16, chapitre~29, décrivent la faillite d'une grande plateforme d'échange). Pour un
salarié, la faillite transforme la rémunération en créance. Aux États-Unis, dans une procédure de faillite,
les salaires, traitements et commissions impayés gagnés dans les 180 jours précédant la requête ne bénéficient d'une priorité parmi les créances non garanties
que jusqu'à un montant fixe par personne, soit 17\,150~\$ depuis le 1\up{er} avril 2025~; le reste est classé avec les autres créanciers non garantis.
Les jetons non acquis et les titres de capital de l'entreprise défaillante ne valent en général rien. Les enseignements pratiques sont ceux de toute
jeune entreprise, en plus aigus~: prendre le numéraire lorsque c'est possible, vendre les jetons acquis selon le calendrier prévu et savoir quelle entité
l'emploie.

\section{Tutoriel~: effectifs et attribution de jetons}\label{tut:in:crypto-native-firms:tut}

\begin{tutorial}
\textbf{Objectif.} Lire les effectifs d'une plateforme d'échange de crypto-actifs au regard de son chiffre d'affaires, et évaluer une attribution de jetons.
\textbf{État final~:} \cref{fig:in:crypto-native-firms:cycle,fig:in:crypto-native-firms:grant}.
\begin{tutsteps}
\item \textbf{Les dépôts.} \texttt{data/industry/crypto\_headcount.csv} contient les effectifs en fin d'année et le chiffre d'affaires total issus
des formulaires 10-K de la plateforme d'échange~; \texttt{in\_crypto.changes()} calcule les variations annuelles.
\item \textbf{L'attribution.} \texttt{firm.tokencomp.Grant(400\_000, months=48, cliff=12, lockup=6, tax=0.40)} et
\texttt{simulate(grant, vol, drift, n, rng)} (\cref{lst:in:crypto-native-firms:sim}).
\omcode{code/firm/tokencomp/firm_tokencomp.py}{39}{58}{Chaque tranche imposée à l'acquisition et vendue à la fin de son blocage, sur
des trajectoires de prix simulées.}\label{lst:in:crypto-native-firms:sim}
\item \textbf{Résumer.} \texttt{firm.tokencomp.summary} donne la médiane, les 10\up{e} et 90\up{e}~centiles, la moyenne,
l'équivalent en numéraire et la part des trajectoires avec une tranche dont l'impôt dépasse son produit.
\end{tutsteps}
\end{tutorial}

\textbf{Pour aller plus loin.} Allonger le blocage à douze mois et observer à quelle fréquence l'impôt dépasse le produit
(exercice~7)~; donner au jeton une dérive positive et chercher quelle dérive rend la médiane égale à la valeur en numéraire.

\section{Construction~: le modèle d'attribution de jetons}\label{bld:in:crypto-native-firms:tokencomp}

\begin{build}
\bfield{Rôle} Donner un chiffre, et une distribution, à la \omterm{def:in:crypto-native-firms:tokens}{rémunération en jetons}, pour ce chapitre et pour la comparaison des
packages de rémunération du chapitre~13.
\bfield{Interface} \texttt{firm.tokencomp}~: \texttt{Grant(value0, months, cliff, lockup, tax, discount)}~;
\texttt{schedule(grant)}~; \texttt{simulate(grant, vol, drift, n, rng)}~; \texttt{summary(sim)}.
\bfield{Règles} L'impôt est dû à l'acquisition, sur la valeur à ce moment-là~; une tranche est vendue à la fin de son blocage~; la dérive et la volatilité du modèle de prix
sont les hypothèses énoncées par l'appelant~; le calendrier totalise l'attribution.
\bfield{Tests d'acceptation} \texttt{code/firm/tokencomp/tests/}~: le calendrier somme à un avec un quart à la carence~; à
volatilité nulle, le résultat égale la valeur en numéraire nette d'impôt~; la volatilité disperse le résultat, abaisse la médiane sous
la valeur en numéraire et maintient la moyenne proche de celle-ci.
\bfield{Pour aller plus loin} Sauts de cours (livre~4)~; un régime fiscal qui impose à la vente~; couvrir une attribution avec des contrats à terme ou des options
là où elles existent et sont autorisées.
\end{build}

\begin{omsources}
\item Coinbase Global, Inc., formulaires 10-K pour 2022, 2023, 2024 et 2025.
\item US Internal Revenue Service, Notice 2014-21~; HM Revenue \& Customs, Cryptoassets Manual, CRYPTO21100.
\item Ethereum Foundation, page d'accueil et EF Report 2024.
\item US Code, titre 11, section 507, et ajustement des montants en dollars par la Judicial Conference, en vigueur le 1\up{er} avril 2025.
\end{omsources}

\section{Exercices}

\begin{exercise}[$\star$]\label{exo:in:crypto-native-firms:1}
Calculer la variation du chiffre d'affaires et des effectifs de la plateforme d'échange pour chaque année de 2022 à 2025.
\end{exercise}

\begin{exercise}[$\star$]\label{exo:in:crypto-native-firms:2}
Dans l'exemple, quelle part de l'attribution est acquise à la carence, et quelle part chaque mois ensuite~?
\end{exercise}

\begin{exercise}[$\star$]\label{exo:in:crypto-native-firms:3}
Un employé reçoit des jetons valant 10\,000~\$ au moment de l'acquisition, imposés à 40~\%, et les vend après une période de blocage au demi-prix.
Quel est le résultat net~?
\end{exercise}

\begin{exercise}[$\star\star$]\label{exo:in:crypto-native-firms:4}
Pourquoi la valeur médiane de l'attribution est-elle inférieure à la valeur en espèces, alors que la moyenne en est proche~?
\end{exercise}

\begin{exercise}[$\star\star$]\label{exo:in:crypto-native-firms:5}
Une entreprise défaillante doit à un employé 60\,000~\$ de salaires gagnés au cours des quatre derniers mois. Quelle part bénéficie d'un privilège selon la règle américaine~?
\end{exercise}

\begin{exercise}[$\star\star$]\label{exo:in:crypto-native-firms:6}
D'après la \cref{fig:in:crypto-native-firms:cycle}, pour quelles deux années les effectifs se sont le plus écartés
du chiffre d'affaires, et qu'est-ce qui explique l'un et l'autre~?
\end{exercise}

\begin{exercise}[$\star\star\star$]\label{exo:in:crypto-native-firms:7}
\emph{Programmation.} Relancer l'attribution avec une période de blocage de douze mois. Comment évoluent la médiane et la part des trajectoires dans lesquelles
l'impôt d'une tranche dépasse son produit~?
\end{exercise}

\begin{exercise}[$\star\star\star$]\label{exo:in:crypto-native-firms:8}
\emph{Trouver la faille.} «~La valeur espérée de l'attribution en jetons est égale à l'offre en espèces, donc les deux offres sont équivalentes.~»
\end{exercise}

\section{Problème~: Rémunération en jetons}

\begin{problem}[{Problème du week-end --- rémunération en jetons}]\label{pb:in:crypto-native-firms:1}
Une ingénieure se voit proposer un salaire plus une attribution de jetons sur quatre ans valant 400\,000~\$ au cours actuel, ou le même salaire plus
400\,000~\$ en espèces sur quatre ans. Elle veut comprendre la différence.

\medskip\noindent\textbf{Partie I --- Les employeurs.}
\begin{enumerate}
\item Nommer les quatre types d'employeurs issus de la crypto.
\item Définir une \omterm{def:in:crypto-native-firms:foundation}{fondation de protocole} et décrire comment une fondation se finance au fil du cycle.
\item Donner quatre façons dont les employeurs issus de la crypto diffèrent du reste de l'industrie.
\item Donner le chiffre d'affaires et les effectifs de la plateforme de 2021 à 2025.
\item Quelles furent les deux restructurations, en proportion des effectifs et en nombre de salariés~?
\end{enumerate}

\medskip\noindent\textbf{Partie II --- La fiscalité.}
\begin{enumerate}[resume]
\item Définir la \omterm{def:in:crypto-native-firms:tokens}{rémunération en jetons} et le \omterm{def:in:crypto-native-firms:tokens}{bon de souscription de jetons} (token warrant).
\item Comment les administrations fiscales américaine et britannique traitent-elles les jetons reçus en contrepartie d'un travail~?
\item Quand l'impôt est-il dû, et sur quelle valeur~?
\item Que se passe-t-il si le jeton baisse pendant la période de blocage~?
\item Calculer le résultat net d'une tranche de 10\,000~\$ imposée à 40~\% et vendue à la moitié de son cours à l'acquisition.
\end{enumerate}

\medskip\noindent\textbf{Partie III --- L'attribution.}
\begin{enumerate}[resume]
\item Écrire le calendrier d'acquisition.
\item Énoncer la méthode d'évaluation de l'attribution.
\item Donner la médiane, les 10\up{e}~et 90\up{e}~centiles et la moyenne à une volatilité de 80~\%.
\item Quel est l'équivalent en espèces, et dans quelle part des trajectoires l'attribution vaut-elle moins~?
\item Dans quelle part des trajectoires l'impôt d'au moins une tranche dépasse-t-il son produit, à une volatilité de 80~\% et à 40~\%~?
\end{enumerate}

\medskip\noindent\textbf{Partie IV --- Le verdict.}
\begin{enumerate}[resume]
\item Énoncer le \emph{résultat nommé}~: la médiane et la fourchette des 10\up{e}--90\up{e}~centiles de l'attribution à une volatilité de 80~\%,
face aux espèces, et la part des trajectoires dans lesquelles l'impôt dépasse le produit.
\item Quelle part de ses salaires impayés bénéficierait d'un privilège si l'entreprise faisait défaut~?
\item Que suggère l'historique des effectifs quant au moment où rejoindre l'entreprise~?
\item Citer deux moyens de réduire le risque de l'attribution, si c'est permis.
\item En deux phrases, comment doit-elle choisir~?
\end{enumerate}
\end{problem}

\section{Questions d'entretien}

\begin{interviewq}[$\star$ \iqroles{développeur}]\label{iq:in:crypto-native-firms:1}
Que fait l'équipe d'ingénierie d'une plateforme de crypto-actifs que celle d'une bourse d'actions ne fait pas~?
\end{interviewq}

\begin{interviewq}[$\star$ \iqroles{trader}]\label{iq:in:crypto-native-firms:2}
Pourquoi les sociétés de trading sur crypto-actifs fonctionnent-elles 24~heures sur 24, et comment organiseriez-vous un desk pour cela~?
\end{interviewq}

\begin{interviewq}[$\star\star$ \iqroles{chercheur, risque}]\label{iq:in:crypto-native-firms:3}
Le prix d'un jeton suit une loi lognormale de volatilité 80~\% et sans dérive. Quelle est la médiane de son prix dans un an, rapporté
à aujourd'hui~?
\end{interviewq}

\begin{interviewq}[$\star\star$ \iqroles{risque}]\label{iq:in:crypto-native-firms:4}
Vous êtes rémunéré en un jeton que vous ne pouvez pas vendre pendant un an. Comment pourriez-vous le couvrir, et qu'est-ce qui vous en empêcherait~?
\end{interviewq}

\begin{interviewq}[$\star\star$ \iqroles{développeur, risque}]\label{iq:in:crypto-native-firms:5}
Une entreprise conserve la majeure partie des actifs de ses clients dans des portefeuilles froids et une partie dans des portefeuilles chauds. Qu'est-ce qui détermine la répartition, et que peut-il mal tourner~?
\end{interviewq}

\begin{interviewq}[$\star\star\star$ \iqroles{chercheur}]\label{iq:in:crypto-native-firms:6}
Avec cinq ans de chiffre d'affaires et d'effectifs d'une même entreprise, comment testeriez-vous si les effectifs suivent avec retard le chiffre d'affaires, et
que concluriez-vous~?
\end{interviewq}
